
We have presented a methodology for testing whether calibration quality mediates the correlation between LLM factuality and adversarial robustness. The approach connects three previously isolated research communities---factuality evaluation, adversarial robustness, and calibration---through a testable hypothesis with quantified success criteria.

The evaluation pipeline is validated and functional. Real TruthfulQA MC1 evaluations were completed for three models:
\begin{itemize}
\item FLAN-T5 base: 0.180
\item FLAN-T5 large: 0.190
\item Phi-2: 0.290
\end{itemize}

However, the core hypothesis remains untested because:
\begin{enumerate}
\item ASR values were simulated rather than measured
\item Only 3 of 12 models were evaluated
\item Mediation analysis was not executed
\end{enumerate}

\textbf{Contribution}: This work provides a validated experimental framework for connecting factuality, robustness, and calibration metrics. The methodology---including correlation analysis with confound control, mediation testing, and intervention validation---can be executed when computational resources permit real TextFooler evaluation.

\textbf{Future Work}: If calibration does mediate the factuality-robustness relationship, this would provide a unified framework for LLM reliability evaluation and suggest that improving calibration could simultaneously improve both dimensions. Completing this evaluation requires executing real adversarial attacks across the full model set.

